\scanpage{93}
\mspage{46}
où $I_{s'}$ \deleted{est le groupe} \inserted{$\simeq\mathcal E_{\widetilde K}$}, $\widetilde{\widetilde K}$ étant le corps
des fractions du hensélisé (\deleted{\illegible{2}} strict)
de $\underline{\mathcal O}_{S,s}$. ($N_{s'}$ \deleted{est celui du \unsure{corps de|complété}} \inserted{s'identifie de $=$ à}
$\mathcal E_{\widetilde K}$, où $\widetilde K$ correspond au hensélisé pas strict)

Ceci posé, on trouve comme précédemment,
si le composé \deleted{de} $\mathcal E_K\mathrel{\text{---}}\mathcal E_{U_{\unsure{\eta|K}}}\mathrel{\text{---}}\mathcal E_U$
s'annule sur $I_{s'}$ \inserted{(a priori, il envoie $I_{s'}$ dans $\pi$)} \deleted{\unsure{pour}} (ce qui signifie
justement \unsure{\deleted{que}|que} $s'\in S'\ldots$) que un hom
$\mathcal E_{k(s)}\longrightarrow\mathcal E_U$, ou si on préfère
$\mathcal E_{k(s)}\longrightarrow\mathcal E_{U_{\widetilde{\mathcal O}_s}}\xleftarrow{\sim}\mathcal E_{X_s}$, qui \deleted{peut} \inserted{donc} s'inter%
\unsure{prète|préter} comme une section de $\mathcal E_{X_s}$ sur
$\mathcal E_{k(s)}$. Et que se passe-t-il pour
$s\notin S'$ ? \deleted{\unsure{Alors}} On a vu que
$D=S\setminus S'$ est purement de codim $1$
(c'est l'image \unsure{réciproque|inverse} du diviseur
$T$ sur $X$ par la fonction $f$ de $X$ sur $S$,
définie si $g\geq1$ ou si $S$ de dim $1\ldots$)
en fait le $\mu=\mu_\alpha$ de tantot, relatif à
\margin{92}
